\documentclass[11pt,letterpaper]{article}
\usepackage[spanish]{babel}
\usepackage[T1]{fontenc}
\usepackage[utf8]{inputenc}
\usepackage[margin=2.3cm]{geometry}
\usepackage{xcolor,booktabs,longtable,enumitem,hyperref}
\definecolor{bosque}{HTML}{174A3A}
\definecolor{coral}{HTML}{D96B4B}
\hypersetup{colorlinks=true,linkcolor=bosque,urlcolor=coral,pdftitle={COBRIA Ecosistema 1.0: especificación de alcance},pdfauthor={Billy Jak Cordero Porras},pdfsubject={Especificación normativa del ecosistema COBRIA},pdfkeywords={COBRIA, arquitectura de software, patrones, NoSQL, trazabilidad}}
\setlist{nosep,leftmargin=*}
\title{\textcolor{bosque}{COBRIA Ecosistema 1.0}\\\large Especificación de alcance, autoridad y conformidad}
\author{Billy Jak Cordero Porras}
\date{Coto Brus, Costa Rica · 2026}
\begin{document}
\maketitle
\begin{abstract}
Este documento identifica las partes normativas del Ecosistema COBRIA 1.0 y sus
relaciones con los productos pedagógicos y científicos. COBRIA es una propuesta
independiente que compone prácticas conocidas mediante contratos verificables; no se
presenta como estándar oficial ni como invención aislada de cada mecanismo.
\end{abstract}
\tableofcontents
\section{Autoridad documental}
El orden de autoridad es: especificación, implementación, evidencia y productos
editoriales. Una explicación del libro o del sitio no puede ampliar un requisito del
núcleo normativo. Una cifra experimental solo puede publicarse si conserva fuente,
configuración, estado y limitaciones.
\begin{longtable}{p{.19\linewidth}p{.27\linewidth}p{.46\linewidth}}
\toprule Producto & Autoridad & Responsabilidad \\ \midrule
Especificación & Normativa & Define requisitos, invariantes, niveles y vocabulario.\\
Implementación & Ejecutable & Materializa contratos sin redefinirlos.\\
Evidencia & Empírica & Registra pruebas, mediciones, fallos y amenazas.\\
Artículo & Científica & Resume contribución, método, resultados y límites.\\
Documento maestro & Metodológica & Desarrolla fundamentos, protocolo y trazabilidad.\\
Libro y sitio & Pedagógica & Enseñan mediante patrones, ejemplos y laboratorios.\\
\bottomrule
\end{longtable}
\section{Invariantes transversales}
\begin{enumerate}
\item El dominio y la aplicación no dependen de interfaz, SDK o almacenamiento.
\item Las claves físicas compactas se traducen exclusivamente en Daters o adaptadores.
\item Toda operación sensible conserva actor, capacidad, ámbito y finalidad.
\item Ningún estado ejecutado se afirma sin un artefacto verificable.
\item Índices, cachés, vectores y proyecciones no adquieren autoridad silenciosamente.
\item Todo derivado declara fuente, contrato, versión, frescura y retiro.
\item La optimización se justifica con mediciones y una alternativa más simple.
\item Las actividades humanas o externas permanecen pendientes hasta recibir evidencia válida.
\end{enumerate}
\section{Módulos del ecosistema}
El ecosistema organiza diecisiete módulos: Manifiesto, Fundamentos, Layers, Pattern
Design, Data, Core, API \& UX, Security, Quality, Performance, Cloud \& Operations,
Analytics \& AI, Engineering, Build, Toolkit, Evidence y Knowledge. No todos son
obligatorios para todos los proyectos; el nivel mínimo debe corresponder al riesgo.
\section{COBRIA Core 1.0}
Core 1.0 permanece congelado. Su ámbito son sistemas NoSQL documentales con datos
canónicos y representaciones reconstruibles. La conformidad se declara por nivel:
\begin{description}
\item[C1 --- Fundamental.] Autoridad, identidad, ámbito, revisión y frontera de escritura.
\item[C2 --- Reconstruible.] C1 más contrato de transformación, equivalencia,
idempotencia, reconstrucción y publicación segura.
\item[C3 --- Gobernado.] C2 más procedencia, frescura, observabilidad, responsable y retiro.
\end{description}
Una implementación solo declara un nivel cuando supera todas sus pruebas obligatorias.
\section{Patrones, capas y datos}
Una ficha de patrón DEBE separar contexto, problema, fuerzas, solución, participantes,
consecuencias, implementación, casos negativos, métricas y relaciones. Las capas DEBEN
respetar dependencias hacia el dominio. El modelado documental DEBE partir de identidad,
ciclo de vida, consultas, invariantes, ámbito y evolución; una estructura de carpetas no
constituye por sí sola una arquitectura.
\section{Seguridad, operación e inteligencia artificial}
Los controles de seguridad se aplican en la frontera y se prueban con casos negativos.
Los ambientes promueven artefactos inmutables y conservan rollback. La telemetría no
debe exponer secretos ni bloquear el flujo principal. Un sistema de IA conserva
procedencia, utiliza capacidades mínimas y se abstiene cuando la evidencia es insuficiente.
\section{Evidencia y estados}
Los estados permitidos son: propuesto, documentado, implementado, ejecutado, reproducido
externamente y revisado por terceros. Ningún estado implica automáticamente el siguiente.
Los resultados locales no demuestran comportamiento distribuido ni superioridad universal.
\section{Productos y versiones}
El libro principal contiene enseñanza duradera; el cuaderno reúne ejercicios; el anexo
técnico conserva trazabilidad de fases. El artículo y el documento maestro consumen las
mismas afirmaciones canónicas. La descarga pública no debe ofrecer versiones finales
competidoras.
\section{Licencia y cita}
Cada tipo de material se rige por la licencia declarada en el repositorio. La atribución
recomendada es: Cordero Porras, B. J. (2026). \emph{COBRIA Ecosistema 1.0: especificación
de alcance, autoridad y conformidad}. Publicación independiente.
\section{Límites}
COBRIA no garantiza rendimiento, seguridad, disponibilidad, cumplimiento legal ni éxito
de un proyecto. No reemplaza análisis de dominio, pruebas en infraestructura real,
consentimiento, evaluación ética o revisión profesional independiente.
\end{document}
